% Stage 8 restored fragment — 2.6.4 Fredholm structure and pattern solvability (archive lines 8332-8619)
\subsection{Fredholm structure and pattern solvability}
\label{sec:ch263}
Pattern synthesis and source inversion are the inverse problems every antenna laboratory
eventually confronts. One specifies a desired far-zone amplitude distribution and asks which
feed currents will produce it; or one measures near fields on a scanner and asks which sources
could have generated them \cite{harrington2001}. When these inversions fail, the
failure is rarely mysterious in hindsight: non-radiating components lie in
$\ker\mathcal{R}_\omega$---the set of admissible impressed currents annihilated by export---and
leave no far-zone signature, so no finite measurement set can resolve them without additional
excitation or constraints \cite{devaney2004}. Recoverable rank from
Subsection~\ref{sec:ch26.signature-aperture} caps the truncated modal system even when a
continuous target is mathematically consistent \cite{hansen1988}. Fredholm theory names
kernel, cokernel, and index after the eigenbasis and rotational structure above are in place
\cite{stratton1941}.

Assumptions. Fix \(\omega>0\), bounded or exterior problem class with outgoing radiation condition on open problems,
physically realizable \(\mathcal{J}_\Omega\) Eq.~\eqref{eq:ch2.source-space-refined}, export map
Eq.~\eqref{eq:ch2.spectral-operator-targets}, and the radiation eigenbasis developed above
\cite{stratton1941}. Source pairing uses dual space
Eq.~\eqref{eq:ch2.dual-source-pairing}; field pairing uses
Eq.~\eqref{eq:ch2.maxwell-adjoint-pairing} \cite{stratton1941}. Compact perturbations are
taken in operator norm on \(\mathcal{J}_\Omega\) for sources with compact support
\(\Omega_s\) \cite{harrington2001}. Regime qualification Eq.~\eqref{eq:ch2.regime-operator-screen}
governs whether recovered \(\mathbf{g}\) may be reported as transport data
\cite{stratton1941}.
The \emph{radiation kernel} $\ker\mathcal{R}_\omega$ collects impressed currents invisible to
qualified far-zone meters. The \emph{cokernel}
$\mathrm{coker}\,\mathcal{R}_\omega=\mathcal{F}_{\mathrm{rad}}/\operatorname{ran}\mathcal{R}_\omega$
labels far-zone pattern targets that no admissible source can realize---equivalently, directions
orthogonal to every radiative image under the adjoint pairing \cite{stratton1941}. The
\emph{Fredholm index}
$\mathrm{ind}\,\mathcal{R}_\omega=\dim\ker\mathcal{R}_\omega-\dim\mathrm{coker}\,\mathcal{R}_\omega$
counts the net imbalance between invisible source degrees and unreachable pattern obstructions;
it is stable under the compact perturbations recorded below \cite{harrington2001}. In the
laboratory, kernel membership explains why two distinct feeds can match the same far pattern;
cokernel membership explains why a synthesis target refuses to converge no matter how many ports
are tuned \cite{stratton1941}.
Source inversion is difficult because radiation smooths fine source detail: the far-zone map
retains only exportable content and discards non-radiating degrees in $\ker\mathcal{R}_\omega$
\cite{devaney2004}. The export map inherits the Green-function split of Section~\ref{sec:ch23}: a leading outgoing
propagation part $\mathcal{R}_{0,\omega}$ plus a compact smoothing part $K_{\mathrm{rad}}$ at
finite separation on bounded source classes \cite{stratton1941}. Physically,
$\mathcal{R}_{0,\omega}$ is the dominant far-zone propagation law and $K_{\mathrm{rad}}$ collects
the compact corrections that smooth fine source structure into resolvable pattern content
\cite{harrington2001}. That split is the operator-side restatement of the range observation that
far-zone data smooth fine source detail. Formally,
\begin{equation}
\label{eq:ch2.radiation-fredholm-decomposition}
\mathcal{R}_\omega
=
\mathcal{R}_{0,\omega}+K_{\mathrm{rad}},
\qquad
K_{\mathrm{rad}}\ \text{compact on}\ \mathcal{J}_\Omega,
\end{equation}
where \(\mathcal{R}_{0,\omega}\) captures leading outgoing singular propagation and
\(K_{\mathrm{rad}}\) collects compact smoothing at finite separation on bounded source classes
\cite{stratton1941}. Equation~\eqref{eq:ch2.radiation-fredholm-decomposition}
specializes Eq.~\eqref{eq:ch2.fredholm-split} to the radiation image
Eq.~\eqref{eq:ch2.integral-radiation-map} after Sections~\ref{sec:ch23}--\ref{sec:ch25}
\cite{stratton1941}. On exterior domains with compact
\(\Omega_s\), \(K_{\mathrm{rad}}\) is compact while \(\mathcal{R}_{0,\omega}\) carries the
non-compact far-zone propagation; on bounded cavities, both parts reduce to discrete modal
books Eq.~\eqref{eq:ch2.cavity-discrete-radiation-spectrum} \cite{harrington2001}.
The Fredholm index counts the net imbalance between invisible source degrees
($\ker\mathcal{R}_\omega$) and unreachable pattern targets
($\mathrm{coker}\,\mathcal{R}_\omega$). On the range, kernel membership means different
feeds may produce identical far patterns; cokernel membership means a synthesis target lies
outside the reachable subspace no matter how many ports are tuned \cite{stratton1941}.
For closed operators with finite-dimensional kernel and cokernel, define the Fredholm index
\begin{equation}
\label{eq:ch2.radiation-fredholm-index-def}
\mathrm{ind}\,\mathcal{R}_\omega
=
\dim\ker\mathcal{R}_\omega
-
\dim\mathrm{coker}\,\mathcal{R}_\omega,
\end{equation}
with \(\mathrm{coker}\,\mathcal{R}_\omega=\mathcal{F}_{\mathrm{rad}}/\operatorname{ran}\mathcal{R}_\omega\)
on the declared quotient by outgoing radiation condition \cite{stratton1941}. Equation
\eqref{eq:ch2.radiation-fredholm-index-def} counts independent non-radiating source degrees
Eq.~\eqref{eq:ch2.ker-R-omega} minus obstructions to reaching a target far pattern
\(\mathbf{g}\in\mathcal{F}_{\mathrm{rad}}\) \cite{stratton1941}. Index is stable under compact
perturbations in Eq.~\eqref{eq:ch2.radiation-fredholm-decomposition}
\cite{harrington2001}. Observation rank cap Eq.~\eqref{eq:ch2.reconstruction-rank-cap}
is independent of \(\mathrm{ind}\,\mathcal{R}_\omega\) but composes with it in inversion
budgets \cite{stratton1941}.
\begin{figure}[t]
\centering
\ChOneFigBlock{%
\begin{tikzpicture}[ch1fig, node distance=5mm and 7mm]
 \node[ch1box] (j) {$\mathcal{J}_\Omega$};
 \node[ch1box, right=7mm of j] (r) {$\mathcal{R}_\omega$};
 \node[ch1box, right=7mm of r] (f) {$\mathcal{F}_{\mathrm{rad}}$};
 \node[ch1box, below=5mm of r] (k) {$K_{\mathrm{rad}}$};
 \node[ch1box, below=5mm of f] (c) {$\mathrm{coker}$};
 \draw[ch1flow] (j) -- (r) -- (f);
 \draw[ch1flow] (k.south) -- (k.south |- r.north);
 \draw[ch1flow] (k.south |- r.north) -- (r.north);
 \draw[ch1flow] (c.south) -- (c.south |- f.north);
 \draw[ch1flow] (c.south |- f.north) -- (f.north);
\end{tikzpicture}%
}
\caption{Fredholm measurement sequence: compact perturbation
Eq.~\eqref{eq:ch2.radiation-fredholm-decomposition}, kernel
\(\ker\mathcal{R}_\omega\), and cokernel obstructions to solvability
Eq.~\eqref{eq:ch2.radiation-solvability-criterion}.}
\label{fig:ch2.fredholm-radiation-measurement chain}
\end{figure}
~\eqref{eq:ch2.rank-composition-law} \cite{stratton1941}.
The Fredholm alternative states that a target far-zone field is synthesizable only if it
carries no component along cokernel directions invisible to the adjoint export map
\cite{stratton1941}. Consequence: a measured pattern residual that is orthogonal to every
reachable image is not a calibration fault---it signals membership in
$\mathrm{coker}\,\mathcal{R}_\omega$ or in a geometrically forbidden sector from
Subsection~\ref{sec:ch264} \cite{harrington2001}. The forward radiation equation $\mathcal{R}_\omega[\Jimp]=\mathbf{g}$
is solvable iff the target $\mathbf{g}\in\mathcal{F}_{\mathrm{rad}}$ is orthogonal to the
cokernel of the adjoint map on declared pairings:
\begin{equation}
\label{eq:ch2.radiation-solvability-criterion}
\mathcal{R}_\omega[\Jimp]=\mathbf{g}
\ \text{solvable}
\iff
\langle \mathbf{g},\,\mathbf{w}\rangle_{\mathcal{U}}=0
\ \ \forall\ \mathbf{w}\in\ker\mathcal{R}_\omega^{\dagger},
\end{equation}
with \(\mathcal{R}_\omega^{\dagger}\) the adjoint on
\((\mathcal{J}_\Omega,\mathcal{F}_{\mathrm{rad}})\) induced by
Eqs.~\eqref{eq:ch2.dual-source-pairing} and~\eqref{eq:ch2.maxwell-adjoint-pairing}
\cite{stratton1941}. Equation~\eqref{eq:ch2.radiation-solvability-criterion}
develops Eq.~\eqref{eq:ch2.spectral-solvability} and aligns with biorthogonality
Eq.~\eqref{eq:ch2.radiation-biorthogonality} on modal adjoints
\cite{stratton1941}. When \(\mathbf{g}\) arises from far-pattern data, solvability is
further capped by Eq.~\eqref{eq:ch2.reconstruction-rank-cap} even if
Eq.~\eqref{eq:ch2.radiation-solvability-criterion} holds in the continuous model
\cite{harrington2001}.
\phantomsection\label{thm:ch2.radiation-fredholm-structure}
Let \(\mathcal{R}_\omega\) satisfy Eq.~\eqref{eq:ch2.radiation-fredholm-decomposition} on
\(\mathcal{J}_\Omega\) with \(\mathcal{R}_{0,\omega}\) Fredholm of index zero on the outgoing
quotient and \(K_{\mathrm{rad}}\) compact. Then \(\mathcal{R}_\omega\) is Fredholm,
\begin{equation}
\label{eq:ch2.radiation-index-stability}
\mathrm{ind}\,\mathcal{R}_\omega
=
\mathrm{ind}\,\mathcal{R}_{0,\omega},
\end{equation}
and Eq.~\eqref{eq:ch2.radiation-solvability-criterion} holds. For finite-rank observation
\(\mathcal{O}_{\mathrm{meas}}:\mathcal{F}_{\mathrm{rad}}\to\mathbb{C}^{N}\),
\begin{equation}
\label{eq:ch2.observation-fredholm-composition}
\mathcal{O}_{\mathrm{meas}}\mathcal{R}_\omega
\ \text{is Fredholm on}\ \mathcal{J}_\Omega
\ \Longleftrightarrow\
\mathcal{R}_\omega\ \text{Fredholm and}\ 
\mathrm{rank}(\mathcal{O}_{\mathrm{meas}})=N,
\end{equation}
with rank bound Eq.~\eqref{eq:ch2.rank-composition-law}
Compact perturbation stability of Fredholm property and index follows from standard operator
theory applied to Eq.~\eqref{eq:ch2.radiation-fredholm-decomposition}; Subsection
~\ref{sec:ch213} recorded the embedding inputs without computing index here
\cite{stratton1941}. Solvability criterion
Eq.~\eqref{eq:ch2.radiation-solvability-criterion} is the Fredholm alternative for
\(\mathcal{R}_\omega\). Composition with finite-rank \(\mathcal{O}_{\mathrm{meas}}\) preserves
Fredholm class and adjusts index by at most \(N\) on the range side, yielding
Eq.~\eqref{eq:ch2.observation-fredholm-composition} and
Eq.~\eqref{eq:ch2.rank-composition-law} \cite{stratton1941}.
The modal fredholm reduction proceeds as follows: on the orthonormal radiation eigenbasis developed above,
\begin{equation}
\label{eq:ch2.modal-fredholm-reduction}
\mathcal{R}_\omega
=
\sum_{n\in\mathcal{N}_{\mathrm{rad}}}
\lambda_n\,P_n\circ\mathcal{S}_\omega,
\qquad
\lambda_n\neq 0
\ \Longleftrightarrow\
\text{mode }n\ \text{exporting},
\end{equation}
so solvability of \(\mathcal{R}_\omega[\Jimp]=\mathbf{g}\) reduces to matching coefficients
\(c_n[\mathfrak{D}]\) Eq.~\eqref{eq:ch2.source-modal-amplitude} on modes with
\(\lambda_n\neq 0\) \cite{stratton1941}. Modes with \(\lambda_n=0\) lie in
\(\ker\mathcal{R}_\omega\) at the source level Eq.~\eqref{eq:ch2.ker-R-omega}; modes
with \(\mathbf{g}\) unsupported on \(\operatorname{ran}\mathcal{R}_\omega\) generate cokernel
obstructions in Eq.~\eqref{eq:ch2.radiation-solvability-criterion}
\cite{stratton1941}. Truncation to \(\mathcal{R}_\omega^{(N)}\) from
Eq.~\eqref{eq:ch2.spectral-compression-operator} yields approximate solvability with error
Eq.~\eqref{eq:ch2.compression-error-bound} controlled by \(\sigma_{N+1}(K_{\mathrm{rad}})\)
\cite{harrington2001}. \phantomsection\label{prop:ch2.modal-invertibility-equivalence}
Let \(\mathbf{g}\in\operatorname{ran}\mathcal{R}_\omega\) and expand
\(\mathbf{g}=\sum_n g_n\,\mathbf{u}_n\) with \(\{\mathbf{u}_n\}\) from
Subsection~\ref{sec:ch261}. Then Eq.~\eqref{eq:ch2.radiation-solvability-criterion} holds if
and only if \(g_n=0\) whenever \(\lambda_n=0\). A particular solution is determined by
\begin{equation}
\label{eq:ch2.modal-particular-solution}
c_n
=
\frac{g_n}{\lambda_n}
\quad\text{for}\ \lambda_n\neq 0,
\end{equation}
with arbitrary coefficients on \(\ker\mathcal{R}_\omega\) not fixed by far-zone data
\cite{stratton1941}. Equation~\eqref{eq:ch2.modal-particular-solution}
is the spectral form of Eq.~\eqref{eq:ch2.inverse-source-equation} on the radiative range;
Tikhonov and regularized inversions from Subsection~\ref{sec:ch244} damp ill-conditioned
\(\lambda_n\) without altering Fredholm index \cite{stratton1941}.
Substitute Eq.~\eqref{eq:ch2.radiation-map-modal-sum} into
\(\mathcal{R}_\omega[\Jimp]=\mathbf{g}\) and use orthogonality of \(\{\mathbf{u}_n\}\) in
\(\mathcal{U}\) \cite{stratton1941}. Modes with \(\lambda_n=0\) do not appear in the range;
modes with \(\lambda_n\neq 0\) invert mode-wise yielding
Eq.~\eqref{eq:ch2.modal-particular-solution}. Kernel directions remain free because
\(\mathcal{R}_\omega\) is not injective on \(\mathcal{J}_\Omega\)
Eq.~\eqref{eq:ch2.ker-R-omega} \cite{harrington2001}.
Cokernel elements pair with unreachable far-zone content: if \(\mathbf{w}\in\ker\mathcal{R}_\omega^{\dagger}\),
then \(\langle \mathbf{g},\mathbf{w}\rangle_{\mathcal{U}}=0\) is required for solvability
Eq.~\eqref{eq:ch2.radiation-solvability-criterion} \cite{stratton1941}. On exterior domains,
continuous angular spectra Eq.~\eqref{eq:ch2.exterior-continuous-spectral-integral} imply
infinite-dimensional cokernel constraints in principle; finite apertures replace them with
rank-\(N\) projections through Eq.~\eqref{eq:ch2.observation-fredholm-composition}
\cite{harrington2001}. Far-pattern coefficients from the eigenbasis express
\(\mathbf{g}\) through \(\mathbf{f}_n\); solvability requires consistency of measured
coefficients with Eq.~\eqref{eq:ch2.modal-particular-solution} on exporting modes only
\cite{stratton1941}.
On the outgoing branch, solvability of \(\mathcal{L}_\omega[\mathbf{u}_\omega]=\mathbf{s}_\omega\)
and solvability of \(\mathcal{R}_\omega[\Jimp]=\mathbf{g}\) are linked through
Eq.~\eqref{eq:ch2.integral-resolvent-equivalence}: forward field construction is single-valued
on \(\mathcal{F}_{\mathrm{SF}}\) while source inversion remains rank-deficient
\cite{stratton1941}. ~\eqref{eq:ch2.coercivity-quotient} is not identical to
\(\mathrm{ind}\,\mathcal{R}_\omega\); export selection removes non-radiative degrees before
pattern inversion \cite{stratton1941}. The Fredholm analysis above treats
\(\mathcal{R}_\omega\) as the primary object for far-zone inverse problems
\cite{harrington2001}.
Fredholm structure on $\mathcal{R}_\omega$ requires linear media, fixed \(\omega\),
outgoing radiation condition, and compactness of \(K_{\mathrm{rad}}\) on the declared source class
\cite{stratton1941}. Resonant cavities violate invertibility of
\(\mathcal{L}_0\) unless quotiented Eq.~\eqref{eq:ch2.coercivity-quotient}; radiation index
must be computed on the outgoing subspace after Eq.~\eqref{eq:ch2.outgoing-radiation-refinement}
\cite{stratton1941}. Loss regularizes small \(\lambda_n\) but does not remove
\(\ker\mathcal{R}_\omega\) Eq.~\eqref{eq:ch2.ker-R-omega} \cite{harrington2001}.
regime qualification Eq.~\eqref{eq:ch2.regime-component-criterion} remains mandatory before transport claims
on recovered \(\mathbf{g}\) \cite{stratton1941}.
As a worked example, a spherical near-field scanner with $N=50$ probe channels is a concrete inversion budget.
The physical configuration places fifty probes on a sphere of radius $R\approx 2\lambda_0$ around
a compact horn feed; each channel records tangential $\E$ or $\H$ depending on probe design
\cite{hansen1988}. After transformation to far-zone coefficients, the composition
$\mathcal{O}_{\mathrm{meas}}\mathcal{R}_\omega$ has rank at most $50$
Eq.~\eqref{eq:ch2.rank-composition-law} \cite{stratton1941}. The observable behavior is
instructive: the first few reconstructed modes match direct far-field cuts taken on the compact
range, while coefficients beyond the fiftieth singular direction track noise rather than new
beam structure.
Even when Eq.~\eqref{eq:ch2.radiation-solvability-criterion} holds for a continuous target
$\mathbf{g}$, discrete recovery satisfies only a truncated modal system; tail error is bounded
by Eq.~\eqref{eq:ch2.compression-error-bound} with the first neglected singular value of the
measurement composition cited from Subsection~\ref{sec:ch26.signature-aperture}
\cite{harrington2001}. The lesson for laboratory planning is direct: adding probes without
enlarging scan radius or improving SNR cannot recover modes the geometry has already annihilated.
Scanner rank is a hard cap that composes with, but does not replace, Fredholm obstructions from
$\ker\mathcal{R}_\omega$ \cite{stratton1941}.
As a worked example, suppose a synthesis target $\mathbf{g}\neq\mathbf{0}$ specifies a far-zone pattern that is
orthogonal to every exporting mode reachable from the declared source class---for example, a
forbidden parity combination above a ground plane when only even sectors are physically
accessible \cite{collin2001}. No amount of feed optimization will satisfy forward
export: the obstruction appears as inconsistent far-zone coefficients in
Eq.~\eqref{eq:ch2.radiation-solvability-criterion}, not as slow convergence of a numerical
solver \cite{harrington2001}. By contrast, when $\mathbf{g}=\mathbf{0}$, the equation is always
solvable with $\Jimp\in\ker\mathcal{R}_\omega$ \cite{stratton1941}; the balanced-loop example
shows that such sources may be strongly measurable in the near zone yet invisible on $\Gamma_R$.
Mode-wise invertibility on the radiative range detects obstructions through $g_n$ on
$\lambda_n=0$ sectors \cite{stratton1941}. Geometric sector constraints on recoverable spectra
are developed next; complete non-radiating classification belongs to Section~\ref{sec:ch272}
\cite{stratton1941}.
\medskip
\noindent
Subsection~\ref{sec:ch263} closed the Fredholm account: compact perturbation, kernel and
cokernel pairings, solvability criterion, and modal reduction on the radiation eigenbasis.
Kernel membership explains invisible source degrees; cokernel membership explains unreachable
pattern targets; recoverable rank from Subsection~\ref{sec:ch26.signature-aperture} caps what
finite instruments can resolve even when continuous solvability holds \cite{stratton1941}.
